\begin{pastexam}[label={px:ust-112-ee-ds-02}]{臺聯大電機類 112 學年度 第 2 題}
（單選；3 分）

If we insert a number 13 to a max-heap [16, 14, 10, 8, 7, 9, 3, 2, 4, 1] at the end of the array
structure [16, 14, 10, 8, 7, 9, 3, 2, 4, 1, 13] and run the necessary procedure to maintain the
max-heap property, what is the new structure after the max-heap property is restored?

(A)~[16, 14, 10, 8, 7, 9, 3, 2, 4, 1, 13]\\
(B)~[16, 14, 10, 8, 7, 9, 13, 2, 4, 1, 3]\\
(C)~[16, 14, 10, 8, 13, 9, 3, 2, 4, 1, 7]\\
(D)~[16, 14, 10, 13, 7, 9, 3, 2, 4, 1, 8]
\end{pastexam}

\begin{pastsoln}
答案是 (C)。

最大堆積插入：新元素放在陣列尾端，然後往上調整（sift-up），只要比父節點大就和父節點交換
（資結書第~8.1~節、虛擬碼~8.2）。 陣列從 $1$ 開始編號，位置 $i$ 的父節點是 $\lfloor i/2 \rfloor$。

\begin{itemize}
\item $13$ 放在位置 $11$，父節點是位置 $5$ 的 $7$。 $13 > 7$，交換：位置 $5$ 變成 $13$，位置 $11$
變成 $7$。
\item $13$ 現在在位置 $5$，父節點是位置 $2$ 的 $14$。 $13 < 14$，停止。
\end{itemize}
結果是 $[16, 14, 10, 8, 13, 9, 3, 2, 4, 1, 7]$，就是 (C)。 檢查：位置 $5$ 的 $13$ 的子節點是位置
$10$、$11$ 的 $1$、$7$，都比 $13$ 小；它的父節點 $14$ 比它大，堆積性質成立。

逐項檢查：
\begin{itemize}
\item (A)：沒有調整。 $13$ 在位置 $11$，父節點 $7 < 13$，不是最大堆積。
\item (B)：把 $13$ 換到位置 $7$。 位置 $11$ 的父節點是 $5$，不是 $7$，$13$ 只會沿著自己到根的
路徑往上走。 而且位置 $7$ 的 $13$ 比它的父節點（位置 $3$ 的 $10$）大，這個陣列根本不是最大堆積。
\item (D)：把 $13$ 換到位置 $4$，同樣走錯路徑。 位置 $11$ 的 $8$ 比它的父節點（位置 $5$ 的 $7$）
大，也不是最大堆積。
\end{itemize}
\end{pastsoln}
